%----------------------------------------------------------------------------------------------------
% START OF CHAPTER
%----------------------------------------------------------------------------------------------------
\chapter{Feb 21}
\paragraph{Topics}
\textit{ATPG. More about articles. Conversation.}
\section{ATPG: All the Possibilities Game}

\begin{table}[ht]
\centering
    \begin{tabular}{ r l  | r l  | r l }
         \itb{tavolo} & \iti{table} & \itb{casa} & \iti{house}  \\
         \hline
         the   & \iti{il tavolo}      & the   & \iti{la casa}         \\
         a     & \iti{un tavolo}      & a     & \iti{una casa}        \\
         good  & \iti{un buon tavolo} & good  & \iti{una buona casa}  \\
         my  & \iti{il mio tavolo}    & good  & \iti{la mia casa}     \\
         expensive  & \iti{un tavolo caro}    & expensive  & \iti{una casa cara} \\
         his   & \iti{il suo tavolo}  & new   & \iti{la sua casa} 
    \end{tabular}
    \caption{ATPG, Feb 21}
    \label{tab:tATPG-Feb21}
\end{table}

\section{Grammar}
\subsection{Articles - singular nouns}
\begin{enumerate}[nosep]
    \item \itb{il} for masculine nouns. \iti{Il gelato}.
    \item \itb{la} for feminine nouns. \iti{La macchina}.
    \item \itb{l'} for nouns starting with vowels. \iti{L'arancia}, \iti{l'albergo}.
    \item \itb{lo} for masculine nouns starting \iti{z} or \iti{s + consonant}.
\end{enumerate}

%--------------------------------------------------------------------------------
% Manual new page: ugly and not nice LaTeX, but makes the document flow better
% when rendered as PDF.
%--------------------------------------------------------------------------------
\newpage

\subsection{Prepositions}
    A \textbf{preposition} is a word that shows the relation of one noun to another. 
    For example, `of', `of the', `on', `with', or `under'. The most flexible in everyday use is
    `of / of the': \iti{di, del, della}, expressing `possession' and sometimes, `specification'.

\begin{itemize}
    \item Claudio's wine. \\
    \iti{Il vino di Claudio.}
    \item Claudio's brother's wine. \\
    \iti{Il vino del fratello di Claudio.}
    \item Claudio's mother's wine. \\
    \iti{Il vino della mamma di Claudio.}
    \item A bottle of red wine. \\
    \iti{Una bottiglia di vino rosso.}
\end{itemize}

\paragraph{on - \iti{su, sul, sulla}} When we have something `over' something.

\begin{itemize}
    \item The wine is on a table. \\
    \iti{Il vino è su un tavolo.}
    \item The wine is on the table. \\
    \iti{Il vino è sul tavolo.}
    \item The wine is on the car. \\
    \iti{Il vino è sulla macchina.}    
\end{itemize}

\paragraph{in - \iti{in, nel, nella}} When we have something `in' something.

\begin{itemize}
    \item The wallet is in the drawer. \\
    \iti{Il portafoglio è nel cassetto.}
    \item The wallet is in the car. \\
    \iti{Il portafoglio è nella macchina.}
    \item The wine is in his car. \\
    \iti{Il vino è nella sua macchina.}
\end{itemize}

\begin{mdframed}
    If you are using the above prepositions with possessive adjectives, you do not need the article. Essentially the preposition includes it `automatically', e.g. `in his car' is \iti{nella sua macchina} and not \iti{nel la sua macchina} nor \iti{nella la sua macchina}.
\end{mdframed}
\paragraph{for, with, between - \iti{per, con, tra}} 
These prepositions are not combined with articles, and act as one word.
\begin{itemize}
    \item I'm talking with my friend. \\
    \iti{Io sto parlando con il mio amico.}
    \item You are driving between home and work today. \\
    \iti{Tu sta guidando tra casa e lavoro oggi.}
    \item I have to work for money. \\
    \iti{Io devo lavorare per soldi.}
\end{itemize}

\section{Vocabulary}
\subsection{Words}

\begin{table}[ht] %[ht] % place table roughly [h]ere or else at the [t]op of page
\centering
    \begin{tabular}{ r l }
          puppy  & \iti{cucciolo} \\
          cat    & \iti{gatto/a}  \\
          today  & \iti{oggi}  \\
          now    & \iti{adesso}
    \end{tabular}
    \caption{Words to practice (Feb 21)}
    \label{tab:tVocabFeb21}
\end{table}

\section{Conversation}
\begin{enumerate}
    \item Nice to meet you. \\
    \iti{Piacere.}\footnote{Also: \iti{È bello incontrare te} (not common).}
    \item My name is Claudio. \\
    \iti{Il mio nome è Claudio.}
    \item I am Italian / American / Italian-American / half-Italian / English. \\
    \iti{Io sono italiano / americano / italo-americano / mezzo-italiano / inglese.}
    \item I work in Virginia. \\
    \iti{Io lavoro in Virginia.}
    \item I am retired \\
    \iti{Io sono in pensione.}
    \item I am in Italy for work. \\
    \iti{Io sono in Italia per lavoro.}
    \item I live in Virginia. \\
    \iti{Io vivo in Virginia.}
    \item I am on vacation for two weeks. \\
    \iti{Io sono in vacanza per due settimane.}
    \item I have two sons. \\
    \iti{Io ho due figli.}
    \item I have one daughter. \\
    \iti{Io ho una figlia.}
    \item This is my wife. \\
    \iti{Questa è mia moglie.}
    \item This is my son. \\
    \iti{Questo è mio figlio.}
    \item This is my friend. \\
    \iti{Questo è il mio amico / Questa è la mia amica.}
    \item This is my mother's house. \\
    \iti{Questa è la casa di mia mamma.}
    \item I speak Italian. \\
    \iti{Io parlo italiano.}
    \item I am very happy. \\
    \iti{Io sono molto felice.}
    \item Italy is very beautiful. \\
    \iti{L'Italia è molto bella.}
    \item I have family in Italy. \\
    \iti{Io ho la famiglia in Italia.}    
\end{enumerate}

\section{Homework}
\begin{enumerate}
    \item I am very happy because I am eating Claudio's pizza. \\
    \iti{Io sono molto felice perché sto mangiando la pizza di Claudio.}
    
    \item Claudio is not happy because I am eating his pizza. \\
    \iti{Claudio non è felice, perché io sto mangiando la sua pizza.}
     
    \item Your glass of wine is on the table. \\
    \iti{Il tuo bicchiere di vino è sul tavolo.}
    
    \item His bottle of wine is very expensive. \\
    \iti{La sua bottiglia di vino è molto cara.}
    
    \item We are studying Italian with Claudio's book. \\
    \iti{Noi stiamo studiando italiano con il libro di Claudio.}
    
    \item My wallet is in the car. \\
    \iti{Il mio portafoglio è nella macchina.}
    
    \item This bread is very good because it is fresh. \\
    \iti{Questo pane è molto buono perché è fresco.}
    
    \item I am speaking with Claudio. \\
    \iti{Io sto parlando con Claudio.}
    
    \item I am sick. (Use the `to stay' verb.) \\
    \iti{Io sto male.}
    
    \item I am sick. (Use the `to be' verb.) \\
    \iti{Io sono malato.}
\end{enumerate}

%----------------------------------------------------------------------------------------------------
% END OF CHAPTER
%----------------------------------------------------------------------------------------------------
